\section{Motivation} \label{sec:motivation_reengineering_gui}
Die bestehende Benutzeroberfläche (User Interface, UI) \todo{Abkürzung UI in acronyms.tex} dient als Schnittstelle für den ursprünglichen Algorithmus. Sie weist funktionale und strukturelle Defizite auf und ist nicht auf den optimierten Algorithmus ausgelegt. Die identifizierten Probleme umfassen mangelhafte Reaktivität, eingeschränkte Navigation sowie architektonische Designmängel.

Ein wesentliches Defizit ist die fehlende Reaktivität. Während der Ausführung des Algorithmus friert die Benutzeroberfläche ein, wodurch jegliche Interaktion während der Rechenzeit blockiert wird. Zudem ist die Skalierung der Ansichtskomponenten statisch. Die grafische Kartenansicht und die Ergebnistabelle besitzen ein festes Größenverhältnis und lassen sich nicht variabel vom Nutzer anpassen. Das sorgt dafür, dass dem Nutzer keine detailierte Sicht auf beide Anzeigen verwäht bleibt. Die Übersichtlichkeit der Applikation wird durch eine fragmentierte Menüführung beeinträchtigt. Konfigurationsparameter verteilen sich inkonsistent auf eine seitliche Menüleiste und auf separate Untermenüs. Darüber hinaus ist die visuelle Zuordnebarkeit der Ergebnisse eingeschränkt. Gipfelkoordinaten aus der Tabelle müssen manuell auf der Karte lokalisiert werden. Eine interaktive Verknüpfung, wie etwa die Hervorhebung eines Gipfels auf der Karte durch die Selektion des zugehörigen Tabelleneintrags, fehlt.

Auf architektonischer Ebene bestehen Mängel in der Codequalität. Die Trennung von Benutzeroberfläche und Geschäftslogik ist unzureichend umgesetzt. Algorithmus-spezifische Verarbeitungsschritte werden direkt in den UI-Komponenten ausgeführt. Obwohl ein objektorientiertes Framework als technologische Basis dient, folgt die Implementierung nicht konsequent objektorientierten Prinzipien. Diese Faktoren mindern die Wartbarkeit sowie die Erweiterbarkeit des Systems. In Summe motivieren diese Defizite ein vollständiges Reengineering der Applikationsoberfläche.